% !TEX root = ../Presentation.tex
% ---------------------------------------------------------------------------
% Summary and next steps
% ---------------------------------------------------------------------------
\section{Summary}

\begin{frame}{Comparison of the models}
  \centering
  \begin{tabular}{lccc}
    \toprule
     & A & B & C \\
    \midrule
    Sessions & pooled & hierarchical & hierarchical \\
    Distractor & separate groups & separate groups & covariate $\beta_2 x_\text{dist}$ \\
    $\text{elpd}_\text{loo}$ (SE) & $-3344.0$ (64.5) & $-3180.0$ (64.3) & $-3217.0$ (65.7) \\
    difference from B (SE) & $-164.1$ (19.2) & best & $-37.1$ (11.9) \\
    effective parameters $p_\text{loo}$ & 13.4 & 178.7 & 138.6 \\
    PSE shift, left / right ($\mu m$) & $-4.19$ / $-6.07$ & $-4.29$ / $-6.06$ & $-4.51$ / $-5.74$ \\
    \bottomrule
  \end{tabular}

  \vspace{0.4em}
  {\scriptsize\color{muted} PSE shift: bimanual $-$ unimanual, posterior mean; for B and C averaged
  over the 42 sessions. Numbers from Fit\_A, Fit\_B and Fit\_C\_Report.}

  \vfill
  \raggedright
  \begin{itemize}
    \item All three models agree on the PSE shift, to within about 0.3~$\mu m$ for each hand
    \item B predicts best: sessions differ (B $\gg$ A), and the shift is not simply proportional to
          distractor amplitude (B $>$ C)
  \end{itemize}
\end{frame}

\begin{frame}{Conclusions}
  \begin{itemize}
    \item The distractor lowers the PSE: by about 4.2--4.5~$\mu m$ (left) and 5.7--6.1~$\mu m$
          (right), in all three models.
    \item It raises the rate of ``high'' responses to weak stimuli ($\gamma_h$).
    \item The sessions differ; a hierarchical model is needed.
    \item The shift grows with distractor amplitude, but a weak distractor still shifts the PSE
          more than a proportional model allows (right hand).
  \end{itemize}
\end{frame}

\begin{frame}{Next steps}
  \begin{itemize}
    \item Revise Model C: lapse rates per condition, a distractor-presence term,
          possibly a bimanual slope.
    \item Model recovery for B against the revised C.
    \item Further experiments: stimulus levels spaced more closely around the PSE, to pin down
          steep slopes, and levels further out, so the curves reach their asymptotes and the
          lapse rates are identified.
  \end{itemize}
\end{frame}
